\begin{papercard}{robillard2019}
\cf{Stages} \stagetag{A} \stagetag{D} \stagetag{C}. Availability, readability and content of privacy policies (PP) \emph{and} terms of agreement (ToA) of mood/mental-health tracking apps.

\cf{Aim} Assess the content and readability of PPs and ToAs of apps that \emph{track} mental-health variables, to characterise how developers handle user data \pc{robillard2019}{2}{the purpose of this study was to assess}.

\cf{Search} Apple App Store and Google Play, using combinations of ``track'' and ``mood'' and their synonyms (Appendix 1). The same searches were repeated on six phone models (iPhone 5c, 5s, 6; Samsung Note 3, OnePlus One, Samsung GS6) to capture device-dependent results \pc{robillard2019}{2}{To capture search results across a range of devices}. The keyword grid crosses 12 action words (track, measure, monitor, test, scale, check, diagnose, trace, log, diary, journal, record) with 15 condition words (mood, anxiety, depression, mental, emotion, feeling, stress, SAD, melancholy, worry, panic, nervous, distress) \pc{robillard2019}{7}{Mood Mood track Mood measure}. The first 100 results of each search were screened \pc{robillard2019}{2}{The first 100 apps that were returned from each search were evaluated}.

\cf{Eligibility} Inclusion: English; free; requires personal information input \emph{over time}; title/description mentions mood, mental health, mental illness, depression, anxiety or a synonym; general public as target users. Exclusion: a condition other than anxiety/depression; a primary non-mental-health focus; one-time input only, i.e.\ no tracking feature \pc{robillard2019}{2}{Inclusion criteria were that the app}. This is a \emph{functional tracking definition} of the domain. PPs/ToAs were located in the app and in associated online content. Documents that applied only to the developer's website were excluded \pc{robillard2019}{2}{In cases where the PPs and ToAs only applied to the developer}.

\cf{Readability} Mean of Gunning Fog, Flesch--Kincaid and SMOG grade levels from an online calculator, following Sunyaev et al.\ \pc{robillard2019}{2}{Grade-level readability of the PPs and ToAs was assessed}.

\cf{Content analysis} A privacy-focused coding guide was developed from a pilot on a random 10\,\% sample \pc{robillard2019}{2}{was developed based on a pilot analysis of a random sample of}. Unit of analysis: the whole document, with a ``rich coding strategy'' that allows multiple categories per document. One investigator coded everything; a second coder coded 20\,\% of PPs and ToAs. Initial agreement (``reproducibility'') was 94\,\% (PP) and 89\,\% (ToA), and disagreements were discussed to consensus \pc{robillard2019}{2}{The entire sample of PPs and ToAs}. Descriptive statistics only.

\cf{Key findings} 369 unique apps. Free iOS: 18\,\% had a PP and 15\,\% a ToA. Free Android: 4\,\% PP and 3\,\% ToA \pc{robillard2019}{2}{For the free iOS apps, 56 PPs}. Mean grade level 13.8 (PP) and 15.2 (ToA). 39\,\% of PPs claimed implied consent through use \pc{robillard2019}{2}{Thirty-nine percent of PPs stated that}. 68\,\% allowed third-party sharing, while only 10\,\% promised consent before sharing \pc{robillard2019}{2}{of PPs stated that users' information may be shared with third parties}. Disclosure occurred in mergers or acquisitions \pc{robillard2019}{2}{the sale or acquisition of the}. Deletion: 16\,\% offered it, and a further 25\,\% only with caveats \pc{robillard2019}{4}{Sixteen percent of PPs}. Security could not be guaranteed according to 37\,\% \pc{robillard2019}{4}{specified that the developer cannot}. The ToAs often granted developers broad licences to user content \pc{robillard2019}{4}{describe a license granted}.

\cf{Limitations stated} Findings were reported only in the document type where a topic appeared more often. PP and ToA content was not consolidated per app, so information missing in one may sit in the other \pc{robillard2019}{6}{we did not consolidate the information in these documents}.

\cf{Use for this thesis} Supports (a) a functional ``tracking over time'' domain rule for mental health, (b) including terms of service alongside policies, (c) device-repeated searches when store results vary, and (d) explicitly coding caveated statements, e.g.\ deletion ``with caveats'' and ``cannot guarantee security''. These are close to the \emph{protection framing} the thesis studies (RQ3). Do not consolidate per app silently: decide and report whether the unit is the document or the app.
\end{papercard}
